\subsection{Choice of parameters}\label{cc:sec:parameters}
The frequencies and orders are chosen so that the truncation and correction
errors are summable. The gains also make each new wave small at its
activation time. We fix the scales here; Appendix~\ref{cc:app:parameters}
proves Lemma~\ref{cc:uf:tower:lem:slow-orders}.

Fix the smooth profiles used in the construction.
The integer $b\ge1$ is the fixed integer specified in
Subsection~\ref{cc:uf:sec:approximate}.
For integers $s\ge1$ put $K(s)=320b(s+1)$.
The exponent $K$ separates successive frequencies.

The estimates of Sections~\ref{cc:sec:localized}, \ref{cc:sec:exact},
\ref{cc:sec:iteration} and~\ref{cc:uf:sec:tower} and of
Lemma~\ref{cc:uf:trunc:lem:cutoff} contain constants that depend only
on an order of differentiation, the depth of the expansion and the
fixed profiles. We call them the fixed-order constants. Choose a
nondecreasing function $\mathcal B$ on the positive integers such that
$\mathcal B(s)$ bounds the fixed-order constants at all orders and
depths at most $7K(s)$, and
\begin{equation}\label{cc:par:floor}
 \log\mathcal B(s)\ge4K(s+1)^3\qquad(s\ge1).
\end{equation}

Let $N\ge\mathcal B(1)$ be the initial frequency parameter. The integer
$q_j$ specifies the spatial derivative order through which step $j$
makes only a small change to the velocity on earlier time intervals.
We define the frequencies and $q_j$ together, starting from $q_1=1$.
Put
\begin{equation}\label{cc:uf:tower:eq:schedule}
 \begin{gathered}
 K_j=K(q_j),\qquad
 \lambda_0=N,\qquad \lambda_j=\lambda_{j-1}^{K_j}\quad(j\ge1),\qquad
 \Lambda_j=\lambda_j^5\quad(j\ge0),\\
 q_{j+1}=\max\{1\le s\le q_j+1:\ \mathcal B(s)\le\lambda_{j-1}\}\quad(j\ge1),\\
 L_j=4\log\lambda_{j+1},\qquad
 \varepsilon_{j-1}=L_j^{-2}
 \quad(j\ge2).
 \end{gathered}
\end{equation}
Here $\lambda_j$, $j\ge1$, is the frequency of wave $j$, and $L_j$ is the
prescribed logarithmic gain. For $j\ge1$, $\Lambda_j$ is the Fourier cutoff
applied to the velocity constructed at step $j$, and $\Lambda_0$ bounds
the frequencies of the shear from which the construction starts.

The orders $q_j$ must tend to infinity for the limit to be smooth;
however fast $\mathcal B$ grows, it only slows their growth. They must
also grow slowly: the amplification estimates of
Section~\ref{cc:uf:sec:tower} lose fixed powers of the logarithmic gain
$L_{j+1}$, which grows with the orders through the exponents $K_j$, and
these losses must stay below small powers of $\lambda_{j-1}$. The lower
bound \eqref{cc:par:floor} is imposed for this purpose.

\begin{lemma}\label{cc:uf:tower:lem:slow-orders}
The integers $q_j$ satisfy $q_j\le q_{j+1}\le q_j+1$ and tend to
infinity, and
\begin{equation}\label{cc:uf:tower:eq:uniform}
 L_{j+1}\le(\log\lambda_{j-1})^2\qquad(j\ge1).
\end{equation}
\end{lemma}
